\section{Panorama general del curso: la automejora es un ciclo cerrado de sistema}

El curso pasa de los modelos de lenguaje a los agentes: un agente tiene metas, entorno, herramientas, estado y retroalimentación, y completa tareas mediante planeación de múltiples pasos y autocorrección. En las primeras ocho sesiones aparece repetidamente el mismo ciclo cerrado:

$$
\text{Propose}\rightarrow \text{Act/Search}\rightarrow
\text{Verify}\rightarrow \text{Learn or Revise}.
$$

El test-time scaling aumenta el muestreo, la búsqueda y la verificación sin cambiar los parámetros; el train-time scaling escribe de vuelta en los parámetros las trayectorias exitosas y el reward; el open-ended search permite además que el sistema explore soluciones o tareas que no estaban enumeradas de antemano.

\videofigure{figures/lecture-roadmap.jpg}{La sesión final se centra en cuatro direcciones futuras: reasoning diversity multiagente, meta-verification matemática, autogeneración de tareas sin datos externos, e intelligence per watt.}{00:03:58--00:05:08}

\videofigure{figures/improvement-bottlenecks.jpg}{La siguiente generación de automejora debe superar tres tipos de cuellos de botella: diversidad de reasoning, robust verification y la selección estática de tareas de entrenamiento.}{00:05:08--00:07:28}

\begin{importantbox}{Los tres recursos escasos de la automejora}
El primero son trayectorias de reasoning novedosas y correctas, el segundo es un verifier confiable, y el tercero son tareas de entrenamiento situadas en el límite actual de la capacidad. Si falta cualquiera de ellos, el ciclo cerrado se estanca prematuramente o se ve destruido por el reward hacking.
\end{importantbox}

\begin{knowledgebox}{El agente es un sistema conjunto entre el LLM y mecanismos externos}
Los LLM actuales suelen ser insuficientes para asumir por sí solos una meta completa. En la práctica, los sistemas orquestan múltiples modelos, herramientas, buscadores, judges y memoria; la capacidad surge del diseño conjunto de los componentes y las rutas de retroalimentación.
\end{knowledgebox}

\subsection{Resumen de la sección}

El valor de largo plazo del curso no está en el artículo más reciente, sino en un conjunto de primitivas reutilizables: generación, búsqueda, herramientas, verificación, reward, memoria, retroceso y entrenamiento. La sesión final muestra, a través de cuatro líneas de investigación, cómo seguirán evolucionando estas primitivas.
